\documentclass[11pt,a4paper]{article}
\usepackage[margin=2.4cm]{geometry}
\usepackage{amsmath,amssymb}
\usepackage{booktabs}
\usepackage{xcolor}
\usepackage[colorlinks=true,linkcolor=black,urlcolor=blue]{hyperref}
\usepackage{enumitem}
\setlist{itemsep=2pt,topsep=3pt}

\newcommand{\win}{\textbf{\color{teal}}}
\newcommand{\lose}{\color{red}}
\newcommand{\Cu}{C_u}
\newcommand{\Cv}{C_v}
\newcommand{\Ce}{C_e}
\newcommand{\ku}{k_u}

\title{\vspace{-1.5cm}Automatic contribution bounding for grouped DP queries}
\author{}
\date{}

\begin{document}
\maketitle
\vspace{-1.2cm}

\section{Motivation}

\subsection{The gap between a DP library and a DP system}

Differentially private aggregation is solved as an \emph{algorithm} and unsolved as a
\emph{product}. Every mature library --- Google DP, OpenDP, Tumult, PipelineDP --- can answer
\texttt{SELECT g, SUM(v) FROM T GROUP BY g} with a rigorous user-level guarantee. But none of them
can answer it \emph{without a human first supplying constants that depend on the data}:

\begin{itemize}
\item \textbf{How much can one user contribute to one group?} A clipping bound. Set it too low and
you destroy the signal; too high and you drown it in noise.
\item \textbf{How many groups can one user touch?} \texttt{max\_partitions\_contributed}, called
$\Cu$ here. A required input in Google DP, with no default.
\end{itemize}

Our system rewrites ordinary SQL into a private plan automatically. An analyst writes a query; they
do not know --- and should not need to know --- the maximum number of months any customer appears
in. So every one of these constants is a blocker, and the question driving this work is narrow and
practical: \textbf{can they be chosen from the data itself, privately, without spending so much
budget that the answer becomes useless?}

\subsection{Why this is hard, and why the constants matter so much}

Two facts make guessing unsafe. First, the penalty is \emph{severely asymmetric}: across our
queries, setting $\Cu$ too small costs up to $38\times$ in error while setting it too large costs
at most $1.7\times$. Second, no single value works. The worst-case penalty of a fixed choice across
our query set:

\begin{center}\small
\begin{tabular}{lrrrrr}
\toprule
analyst supplies & $\Cu{=}1$ & $\Cu{=}10$ & $\Cu{=}100$ & $\Cu{=}$ schema max & automatic \\
\midrule
worst case over queries & 177.6$\times$ & 136.3$\times$ & 5.18$\times$ & 4.18$\times$ &
\textbf{1.36$\times$} \\
\bottomrule
\end{tabular}
\end{center}

$\Cu{=}1$ is catastrophic on TPC-H and simultaneously \emph{optimal} on ClickBench. A single
constant cannot serve both, so the choice must adapt to the data --- which means spending privacy
budget to make it.

\subsection{Where the budget goes}

A grouped DP query has three channels, and every idea in this document buys one of them at the
expense of another:

\begin{enumerate}
\item \textbf{Bound selection} ($\epsilon_b$): how large is one user's contribution?
\item \textbf{Partition selection} ($\epsilon_\eta$): which groups may be named at all? A group
existing because of a single user must not be revealed, so groups are released only when a noised
count of contributing users clears a threshold $\tau$.
\item \textbf{Values} ($\epsilon_v$): the noised aggregates themselves.
\end{enumerate}

The third is what the analyst wants; the first two are overhead. \textbf{Reducing that overhead, or
removing the need for it, is the whole game.}

\subsection{The baseline}

Google DP as published (Wilson et al., PoPETs 2020): \textsc{ApproxBounds} over per-cell values in
base-2 log bins gives $U$; a single $\Cu$ randomly truncates each user to $\Cu$ groups, applied to
\emph{both} the value and the vote channel; value sensitivity is $\Cu \cdot U$ with Laplace noise;
partition selection uses the Wilson threshold
$\tau = 1 - \Cu\log(2-2(1-\delta_\eta)^{1/\Cu})/\epsilon_\eta$.

We compare against this, not against a hardened variant of our own invention. Where we do build a
stronger baseline (\S\ref{sec:decomp}) we say so explicitly, because a comparison against a
mechanism nobody has published or deployed is not evidence about the state of the art.

\paragraph{Scoring.} Throughout: relative $\ell_1$ error over the \emph{true, uncapped} key set,
with $\tau$-suppressed groups charged as released $0$ (full error). Every arm is tuned separately
over its own free parameters --- $\Cu$ or $(\Ce,\Cv)$, and the three-way budget split. Anything
less lets a mechanism buy accuracy by returning fewer rows, or flatters one side by fixing a
parameter the other gets to optimise.

\section{Four ideas, and which channel each attacks}

Ideas (A) attacks the value channel; (B) and (D) attack partition selection; (C) attacks bound
selection. They are independent, and \S\ref{sec:decomp} shows each helps the baseline as much as it
helps us.

\paragraph{(A) $\ell_1$ per-unit norm clip.} Rather than bounding each (unit, group) cell by $U$
and multiplying sensitivity by $\Cu$, bound the unit's \emph{total} across groups:
\[
n_u = \sum_g \bigl|\mathrm{clip}(t(u,g), -B, B)\bigr|,
\qquad
v(u,g) = \mathrm{clip}(t(u,g),-B,B)\cdot\min\!\bigl(1, B/n_u\bigr),
\]
giving $\|v_u\|_1 \le B$ regardless of how many groups $u$ touches. $B$ comes from
\textsc{ApproxBounds} applied to per-unit \emph{norms} rather than per-cell values.

\paragraph{(B) Frozen partition set.} Release the group set once from the filterless query under
its own $(\epsilon_0,\delta_0)$, then reuse it: later queries skip $\tau$ entirely for those
groups, freeing $\epsilon_\eta$ for the values. Sound because partition selection protects group
\emph{existence}, not values --- once released, $G_{\mathrm{fix}}$ is public and conditioning on
it is post-processing.

\paragraph{(C) EM selection of $\Cu$.} Score each candidate $r$ by $q(d,r) = -|F_d(r) - T|$ where
$F_d(r) = |\{u : \ku \le r\}|$ and $T$ is a target rank, then sample by exponential mechanism.
The key observation: removing one unit shifts $K-k+1$ cumulative counts at once, so
\emph{releasing} $(F(1),\dots,F(K))$ would have $\ell_1$ sensitivity up to $K$ --- but the EM
scores each candidate separately, so $\Delta q = \max_r |q(d,r)-q(d',r)| \le 1$.

\paragraph{(D) All-or-Frozen thresholding.} Replace $K$ per-group release tests with one global
AND: release the query-specific group set iff \emph{every} group clears a threshold, else fall
back to $G_{\mathrm{fix}}$. Conventional thresholding leaks if \emph{any} of a unit's $\le\Cu$ new
groups passes (an OR, needing $\rho_\tau \le \delta/\Cu$); the AND needs only $\rho_\tau\le\delta$,
saving $\approx(\Cu/\epsilon_B)\log \Cu$ on the threshold.

\section{Headline results}

\subsection{Utility versus the published baseline}

Fourteen queries, three datasets (TPC-H sf10, StackOverflow, ClickBench), SUM and COUNT measures,
$\epsilon=1$, $\delta=10^{-6}$. \textsc{g-best} gives the baseline its \emph{optimal} $\Cu$ per
query --- an upper bound on any configuration an analyst could supply.

\begin{center}\small
\begin{tabular}{lrrrrr}
\toprule
query & groups & g-best & ours & ours+frozen & vs g-best \\
\midrule
tpch SUM price / month        &     84 &  0.68\% & 0.17\% & \textbf{0.16\%} & 4.38$\times$ \\
tpch SUM price / mo$|$nation  &  2,095 & 13.96\% & 3.95\% & \textbf{2.25\%} & 6.21$\times$ \\
tpch SUM price / mo$|$prio    &    420 & 10.10\% & 3.11\% & \textbf{1.83\%} & 5.53$\times$ \\
tpch SUM price / day          &  2,526 & 99.56\% & 85.22\% & \textbf{10.41\%} & \textbf{9.56$\times$} \\
tpch COUNT / mo$|$nation      &  2,096 & 14.76\% & 4.39\% & \textbf{2.50\%} & 5.90$\times$ \\
tpch SUM qty / mo$|$nation    &  2,096 & 15.52\% & 5.31\% & \textbf{3.17\%} & 4.90$\times$ \\
tpch SUM price / yr$|$nation  &    175 &  0.70\% & 0.37\% & \textbf{0.33\%} & 2.14$\times$ \\
so COUNT posts / month        &    190 & 41.55\% & 36.36\% & \textbf{29.74\%} & 1.40$\times$ \\
so SUM score / month          &    190 & 62.73\% & 45.81\% & \textbf{37.86\%} & 1.66$\times$ \\
so COUNT comments / month     &    182 & 55.21\% & 48.12\% & \textbf{35.76\%} & 1.54$\times$ \\
so COUNT posts / day          &  5,107 & 92.65\% & 89.44\% & \textbf{56.29\%} & 1.65$\times$ \\
cb COUNT hits / region        &  3,238 &  6.00\% & \textbf{5.72\%} & 7.27\% & 1.05$\times$ \\
cb COUNT hits / date$|$reg    &  7,564 & 10.01\% & \textbf{8.49\%} & 12.38\% & 1.18$\times$ \\
cb SUM width / region         &  3,238 &  6.41\% & \textbf{6.29\%} & 9.50\% & 1.02$\times$ \\
\midrule
\multicolumn{5}{l}{\textbf{median 1.90$\times$, range 1.02--9.56$\times$}} \\
\bottomrule
\end{tabular}
\end{center}

The best row --- \texttt{tpch SUM price / day}, $99.56\% \to 10.41\%$ --- is the regime where
$\tau$ suppresses nearly everything and the frozen set is the only thing that works.

\subsection{The scope condition}

The distribution of $\ku$ (groups touched per unit) predicts every row, and we established it in
three parts.

\begin{center}\small
\begin{tabular}{llll}
\toprule
$\ku$ distribution & example & median gain & winner \\
\midrule
concentrated, high ($\gtrsim 10$) & TPC-H (med 30, max 72) & 5.78$\times$ & ours \\
concentrated, low (all touch one) & ClickBench (med 1, max 54) & 1.06$\times$ & tie \\
bimodal (most one, few many) & constructed; so ViewCount/day & 0.68--0.97$\times$ & \lose baseline \\
\bottomrule
\end{tabular}
\end{center}

Sweeping fifteen $\ku$ families at matched unit count, group count and value distribution isolates
two effects that a one-sided ``concentration'' story missed:

\begin{itemize}
\item \textbf{A level threshold near $\ku \approx 7$.} \texttt{constant(5)} is \emph{perfectly}
concentrated yet the baseline wins $0.79\times$; \texttt{constant(20)} and \texttt{constant(60)}
lose $1.5$--$1.6\times$. The clip needs units touching roughly ten or more groups before it pays.
\item \textbf{A genuine loss region.} With most units at $\ku=1$ plus a $2\%$ tail at $\ku=100$,
the baseline wins $1.27\times$: the whales set our $B$ and every ordinary unit pays that noise,
whereas truncation discards them for the price of $U$. Our clip refuses to discard anyone, which
is here the wrong instinct.
\end{itemize}

\subsection{Decomposition: what the gain is actually made of}\label{sec:decomp}

Two decompositions matter for honest reporting.

\paragraph{Bounds are not free.} $U$ and $B$ both come from \textsc{ApproxBounds}, a DP release
costing $\epsilon_b$. Giving \emph{both} arms oracle bounds on a $1.19\times$-spaced grid:

\begin{center}\small
\begin{tabular}{lrrr}
\toprule
query & DP bounds & oracle bounds \\
\midrule
tpch SUM price / mo$|$nation & 3.57$\times$ & \textbf{2.25$\times$} \\
so COUNT posts / month & 0.94$\times$ & 1.02$\times$ \\
cb COUNT hits / region & 1.08$\times$ & 1.05$\times$ \\
\bottomrule
\end{tabular}
\end{center}

So on TPC-H roughly $2.25\times$ is the clipping geometry, surviving perfect bounds on both sides,
and $\approx 1.6\times$ is \textsc{ApproxBounds} serving the baseline's per-cell $U$ worse than our
per-unit $B$. The latter has a cause: \textsc{ApproxBounds} is a max-finder while the error-optimal
per-cell bound is roughly half the maximum, making it $\approx 2\times$ suboptimal for $U$ by
construction. \emph{Ours is less sensitive to bound selection}, which is a robustness claim, not a
sharper clip.

\paragraph{The ideas are independent of each other.} The frozen set (B) and the EM (C) act on the
\emph{partition-selection} channel; the clip (A) acts on the \emph{value} channel. Neither depends
on the other, and both help the published baseline just as much: giving it the frozen set gives
$4.25\%\to3.52\%$, $13.50\%\to9.76\%$ and $82.10\%\to20.20\%$ across three filters, a
$1.2$--$4.1\times$ gain with no $\ell_1$ clip anywhere. Idea (C) is if anything \emph{more} useful
to the baseline than to us, since our clip has no $\Cv$ to select. They should therefore be
reported --- and adopted --- separately.

\paragraph{Most of the stack's gain is the frozen set, which is prior art.} Against a maximally
strong baseline (top-$\Cv$ selection, rescale-with-re-clip, decoupled $\Ce/\Cv$ --- none of which
is published), the clip alone is a wash: median $1.01\times$. Against the published baseline the
clip alone is $1.23\times$. The larger stack numbers come from the frozen partition set.

\subsection{Sensitivity to $\epsilon$ and to the metric}

Two robustness checks that could have invalidated everything:

\begin{center}\small
\begin{tabular}{lrrrrr}
\toprule
& $\epsilon{=}0.25$ & $0.5$ & $1$ & $2$ & $4$ \\
\midrule
median gain & 1.20$\times$ & 2.11$\times$ & 2.88$\times$ & 2.39$\times$ & 3.08$\times$ \\
\bottomrule
\end{tabular}
\qquad
\begin{tabular}{lr}
\toprule
metric & median \\
\midrule
total-$\ell_1$ & 1.23$\times$ \\
mean per-group rel. & 1.19$\times$ \\
median per-group rel. & 1.22$\times$ \\
nRMSE & 1.29$\times$ \\
\bottomrule
\end{tabular}
\end{center}

The ranking is stable in $\epsilon$ (TPC-H wins everywhere, ClickBench ties everywhere) and the
gain is \emph{smallest} at small $\epsilon$ --- the opposite of a resonance peak. All four metrics
agree in direction, with per-group metrics \emph{more} favourable on TPC-H
($5.17\times$ median-rel against $3.61\times$ total-$\ell_1$).

\section{Automatic $\Cu$ selection}

The EM (idea C) is sound and nearly free, but the target fraction matters. Penalty is error at the
selected $\Cu$ over error at that query's best $\Cu$, so $1.00\times$ equals an oracle:

\begin{center}\small
\begin{tabular}{lrrrrr}
\toprule
query & best $\Cu$ & $p{=}0.7$ & $p{=}0.9$ & $p{=}0.95$ & $p{=}0.99$ \\
\midrule
tpch month           & 55 & \lose 17.75$\times$ & 2.91$\times$ & 2.91$\times$ & \textbf{1.35$\times$} \\
tpch month$|$nation  & 44 & 1.18$\times$ & 1.02$\times$ & \textbf{1.00$\times$} & 1.21$\times$ \\
tpch month$|$prio    & 55 & 1.41$\times$ & 1.03$\times$ & \textbf{1.00$\times$} & 1.17$\times$ \\
so posts$|$month     & 44 & 2.02$\times$ & 1.56$\times$ & 1.31$\times$ & \textbf{1.22$\times$} \\
so comments$|$month  & 34 & 1.70$\times$ & 1.34$\times$ & \textbf{1.11$\times$} & 1.36$\times$ \\
cb hits$|$region     &  1 & 1.02$\times$ & 1.00$\times$ & \textbf{0.99$\times$} & 1.00$\times$ \\
\midrule
\textbf{worst case} & & \textbf{17.75$\times$} & 2.91$\times$ & 2.91$\times$ & \textbf{1.36$\times$} \\
\bottomrule
\end{tabular}
\end{center}

The tolerance band is severely asymmetric --- halving $\Cu$ costs $2.6$--$35\times$, doubling costs
only $1.7$--$2.0\times$ --- so the selection should sit deliberately above the optimum. Note the
choice interacts with the mechanism: if $\Cu$ is shared with $\tau$ (the published baseline),
$p\approx0.95$ is better, since $p=0.99$ chases heavy tails and inflates $\tau$; with $\Ce$ and
$\Cv$ separate, $p\approx0.99$ wins.

\paragraph{Why automate at all.} No fixed value works across datasets. Worst-case penalty over
queries: $\Cu{=}1$ gives $177.6\times$, $\Cu{=}10$ gives $136.3\times$, $\Cu{=}100$ gives
$5.18\times$, $\Cu{=}\text{schema max}$ gives $4.18\times$, and the EM gives $\mathbf{1.36\times}$.
$\Cu{=}1$ is catastrophic on TPC-H yet optimal on ClickBench. The mechanism wins by
\emph{adapting}, not by being cleverer than the analyst.

\paragraph{Frozen versus per-query.} A filter on the privacy unit leaves the $\ku$ distribution
\emph{identical} (p50 $=30$, p95 $=48$ even at a filter retaining $9{,}182$ of $1{,}000{,}000$
units), so re-running the EM returns the same value and is slightly worse for paying its cost
again. A filter that removes \emph{groups} collapses $\ku$ (p50 $30\to4$) and per-query wins
$1.8\times$. The decidable rule: re-run selection when the filter constrains a grouping-key
component \emph{not} functionally determined by the privacy unit.

\section{All-or-Frozen: correct but inert}

The privacy analysis is right --- we reproduced the predicted threshold gap to the digit
($\tau_{PG}-\tau_{AF} = 6{,}158$ measured against $(\Cu/\epsilon_B)\log\Cu = 6{,}158$) --- but the
rule essentially never fires. Across eight queries and three datasets, as specified it fires on
none; with a dedicated distinct-unit count instead of the reused bounding histogram it fires on
one; with additional blocking by a unit-determined key component it reaches $20\%$ of blocks on
one more.

The reason is structural and needs no simulation. Any valid threshold satisfies
$\Pr[1 + \mathrm{Lap} \ge \tau] \le \delta$, so $\tau \gg 1$ and \emph{a group holding a single
privacy unit can never pass}. Five of eight queries have minimum support exactly $1$. But
single-unit groups are precisely why partition selection exists.

\begin{quote}
All-or-Frozen requires a query with no singleton group --- and such a query does not need partition
selection in the first place.
\end{quote}

The one query where it fires is also the one where $\tau$ suppresses nothing, so the two group sets
coincide and the result is vacuous. Two findings are worth keeping: reusing the bounding histogram
costs $736\times$ on the threshold relative to a dedicated count (a free statistic is not free when
it is the wrong statistic), and blocking by a unit-determined key component is a valid way to
shrink any per-query conjunction.

\section{Privacy validation}

Two bugs were found in \emph{our own} mechanism, both by explicit neighbour construction after the
argument had already convinced us.

\paragraph{Bug 1: the clip was not a norm.} Written as $n_u = \sum_g \min(t,B)$ it is a signed sum.
On signed measures it fails open (cancelling cells give $n_u=0$, so the unit is released unclipped)
or diverges ($n_u<0$ gives a negative scale). On \texttt{SUM(shipped - returned)} by (customer,
month) at sf10, \textbf{393,053 of 999,982 units exceeded the bound}, worst $2.6\times$ --- making
the release $2.6\epsilon$-DP. Fixed by $n_u := \sum_g|\mathrm{clip}(t,-B,B)|$, which is
bit-identical on non-negative data.

\paragraph{Bug 2: $\delta$ under-charged, demonstrated by an adversary.} A unit creates a group by
having a \emph{value} there, not a vote, so with votes truncated to $\Ce$ the $\ku-\Ce$ unvoted
groups each get an independent chance to clear $\tau$ on noise alone --- chances $\tau$'s union
bound never covered. An explicit attacker over $4\times10^6$ trials:

\begin{center}\small
\begin{tabular}{rrrrr}
\toprule
$\Ce$ & $\ku$ & leak, no gate & leak, gated & over by \\
\midrule
1 & 30 & $2.10\times10^{-5}$ & $1.00\times10^{-6}$ & \lose 21.0$\times$ \\
1 & 72 & $4.70\times10^{-5}$ & $1.25\times10^{-6}$ & \lose \textbf{47.0$\times$} \\
5 & 72 & $1.27\times10^{-5}$ & $1.75\times10^{-6}$ & \lose 12.8$\times$ \\
\bottomrule
\end{tabular}
\end{center}

The fix is a one-line gate --- release a group only if it has at least one vote --- which is free:
on real data it removes 19 of 2,095 groups, each with a $2\times10^{-7}$ chance of release.

\paragraph{What else the adversary could not break.} Value channel: advantage $0.2587$ against a
$0.2913$ bound, and \emph{scale-free} across a $128\times$ range of $B$. Omniscient attacker
(knows all other units exactly): $0.2590$ even for a target trying to contribute $4B$, because the
clip caps it. Reconstruction: posterior sd $1{,}167$ against prior $1{,}183$, a $1.01\times$
narrowing. Reuse of one frozen set across $N$ queries: advantage tracks the composed $\epsilon_v$
only. Adaptivity, where round 1 selects the target for round 2: \emph{loses} to a blind attacker,
$0.36\times$ at $M{=}2$ down to $0.01\times$ at $M{=}100$. SQL join fan-out, with one customer
owning up to 182 rows: $\max\|\text{released}(D)-\text{released}(D\setminus u)\|_1 = 1.000000\,B$
exactly, because the clip is on the unit's total rather than on rows.

\paragraph{Refreshes must be budgeted.} $R$ re-releases of the frozen set give $R$ independent
chances: $\Pr[\text{exposed}] = 1-(1-p)^R \le R\,\delta_0$, confirmed exactly at $R=1..50$. Batching
updates does not reduce the cost of a single refresh; it reduces $R$.

\section{Methodology}\label{sec:method}

Nine headline numbers were produced and then invalidated during this work:
$880\times \to 12\times \to 4.8\times \to 4.65\times \to 1.36\times$, and separately
$4.5\times \to 1.9\times$. Every collapse had one of three causes, and we state them because the
pattern is more transferable than any of the numbers.

\begin{enumerate}
\item \textbf{An untuned baseline.} Comparisons at a fixed budget split, a fixed $\Cu$, or against
a library default are not evidence. Each free parameter must be optimised \emph{for the baseline}
before the comparison is meaningful --- $\Cu$ tuning alone moved one result by $12\times$.
\item \textbf{A forgiving scoring rule.} Scoring against a capped or truncated ``truth'' silently
forgives exactly the bias being measured. Suppressed groups must be charged full error, or a
mechanism can buy accuracy by returning fewer rows.
\item \textbf{Mixing regimes in an average.} A median across queries that violate the scope
condition misleads in both directions. Report by regime.
\end{enumerate}

Three further practices proved necessary:

\begin{itemize}
\item \textbf{Verify sensitivity by neighbour simulation, never by argument.} Both privacy bugs
survived a convincing derivation and died on explicit construction.
\item \textbf{At $\delta$ scale, compute exactly rather than simulate.} A Monte Carlo run
suggested a $1.4\times$ violation that was pure noise; resolving a $1.4\times$ effect at
$\delta=10^{-6}$ needs $\sim10^8$ trials.
\item \textbf{Check the grid before believing a negative.} An ``oracle bounds'' comparison on a
$2\times$-spaced grid --- coarser than the $1.3\times$ error basin --- reported our arm getting
worse under perfect bounds. Pure artifact.
\end{itemize}

\section{Prior art}

Four of five techniques exist. This was checked against library source and primary papers.

\begin{center}\small
\begin{tabular}{lll}
\toprule
technique & status & where \\
\midrule
frozen partition set & \lose ships & Tumult \texttt{get\_groups}+\texttt{KeySet}; Beam \texttt{SelectPartitions} \\
Gaussian partition selection & \lose published & Google C++, OpenDP, PipelineDP, Qrlew; Desfontaines 2022 \\
adaptive log base & \lose config flag & \texttt{ApproxBounds::SetBase/SetScale/SetNumBins} \\
$\ell_1$ norm clipping & \lose adjacent & Harrison \& Manurangsi, $\ell_r$-constrained contributions \\
count-conditioned shrinkage & unclaimed & nearest: Private-PGM; FORC 2025 debiasing \\
\bottomrule
\end{tabular}
\end{center}

Notably, Desfontaines et al.\ (PoPETs 2022) already publish the Gaussian-versus-Laplace crossover
at $\kappa=3$ with $\epsilon/\delta$-invariance; we independently measured $\ku \approx 4.3$ with
the same invariance, which is a rediscovery. What remains unclaimed is narrower: a
\emph{persistent} key set reused across separately-submitted queries (everyone frames reuse within
one pipeline), the tuning rule for the log base, and the count-shrinkage post-processing.

\section{Recommendations}

\begin{enumerate}
\item \textbf{Adopt EM selection of $\Cu$} with $p\approx0.95$ (shared $\Cu$) or $p\approx0.99$
(separate $\Ce/\Cv$). It costs $\approx1\%$ of $\epsilon$, beats any fixed value by $\approx3\times$
worst-case, and its stability premise holds for the quantity it estimates.
\item \textbf{Adopt the persistent frozen partition set}, with public predicate pruning when the
filter constrains the grouping key, and a guard that skips it entirely for unfiltered queries
(where it \emph{hurts}: $5.71\% \to 7.29\%$ on ClickBench).
\item \textbf{Do not pursue All-or-Frozen.} Correct but inert, for a structural reason.
\item \textbf{Ship the $\ell_1$ clip with both fixes}, conditioned on the $\ku$ scope test, and
present its value as \emph{parameter elimination} --- it removes a bound that is unfreezable,
query-dependent, and worth up to $38\times$ when misspecified --- rather than as raw utility.
\item \textbf{Report by regime, and report the published baseline.} In scope:
$3.3$--$3.6\times$ versus published, $1.3$--$2.3\times$ versus the strongest baseline we could
build. Out of scope: a wash or a small loss.
\end{enumerate}

\paragraph{Open.} Nothing is implemented in \texttt{src/}; both privacy fixes must land
\emph{with} the clip, not after. Untested: the smooth-sensitivity (SASS) path, where a preliminary
port \emph{lost} because its value channel is independently broken; workloads with genuine
seasonal alignment, where the baseline's rescaling assumption should fail and our advantage should
grow; and drift-triggered rather than scheduled refreshes.

\end{document}
